\textbf{Phần 1. Trắc nghiệm lựa chọn}
\setcounterpageref{ex}{0}

\begin{ex} %[2H3N17]
Cặp sách đi học của em cân nặng khoảng
\choice
{$30\text{ kg}$}
{$300\text{ kg}$}
{\True $3\text{ kg}$}
{$300\text{ g}$ (quá nhẹ)}
\loigiai{
Cặp sách học sinh nặng khoảng $3\text{ kg}$.\\
\textbf{Đáp án của câu hỏi là C}
}
\end{ex}

\begin{ex} %[2H3N17]
Một xô nước đầy có dung tích khoảng
\choice
{$1\text{ l}$}
{$100\text{ l}$}
{$1000\text{ l}$}
{\True $10\text{ l}$}
\loigiai{
Xô nước chứa khoảng $10\text{ l}$.\\
\textbf{Đáp án của câu hỏi là D}
}
\end{ex}

\begin{ex} %[2H3H17]
Con gà cân nặng $2\text{ kg}$, con ngỗng nặng hơn con gà $3\text{ kg}$. Con ngỗng cân nặng là
\choice
{$4\text{ kg}$}
{\True $5\text{ kg}$}
{$6\text{ kg}$}
{$7\text{ kg}$}
\loigiai{
$2 + 3 = 5\text{ kg}$.\\
\textbf{Đáp án của câu hỏi là B}
}
\end{ex}

\begin{ex} %[2H3H17]
Một chai nước ngọt chứa $2\text{ l}$, rót đầy vào 4 cốc như nhau. Mỗi cốc chứa được
\choice
{\True Nửa lít ($500\text{ ml}$)}
{$1\text{ l}$}
{$2\text{ l}$}
{$4\text{ l}$}
\loigiai{
Rót đều được nửa lít mỗi cốc.\\
\textbf{Đáp án của câu hỏi là A}
}
\end{ex}

\begin{ex} %[2H3H17]
Một túi cam nặng $3\text{ kg}$, 4 túi cam như thế nặng tất cả là
\choice
{$10\text{ kg}$}
{\True $12\text{ kg}$}
{$14\text{ kg}$}
{$15\text{ kg}$}
\loigiai{
$3 \times 4 = 12\text{ kg}$.\\
\textbf{Đáp án của câu hỏi là B}
}
\end{ex}

\begin{ex} %[2H3H17]
Có 3 can nước, mỗi can chứa $5\text{ l}$. Cả 3 can chứa số lít nước là
\choice
{$12\text{ l}$}
{$18\text{ l}$}
{$20\text{ l}$}
{\True $15\text{ l}$}
\loigiai{
$5 \times 3 = 15\text{ l}$.\\
\textbf{Đáp án của câu hỏi là D}
}
\end{ex}

\begin{ex} %[2H3V17]
Mẹ mua một con cá nặng $2\text{ kg}$ và một miếng thịt nặng $1\text{ kg}$. Tổng khối lượng mẹ đã mua là
\choice
{$2\text{ kg}$}
{$4\text{ kg}$}
{\True $3\text{ kg}$}
{$5\text{ kg}$}
\loigiai{
$2 + 1 = 3\text{ kg}$.\\
\textbf{Đáp án của câu hỏi là C}
}
\end{ex}

\begin{ex} %[2H3V17]
Bình chứa $8\text{ l}$ nước ép táo, rót ra mời khách hết $5\text{ l}$. Bình còn lại số lít nước ép là
\choice
{$2\text{ l}$}
{$4\text{ l}$}
{$5\text{ l}$}
{\True $3\text{ l}$}
\loigiai{
$8 - 5 = 3\text{ l}$.\\
\textbf{Đáp án của câu hỏi là D}
}
\end{ex}

\begin{ex} %[2H3V17]
Trên đĩa cân bên trái có 1 quả dưa hấu, đĩa bên phải có 2 quả cân loại $2\text{ kg}$ thì cân thăng bằng. Quả dưa hấu nặng là
\choice
{$2\text{ kg}$}
{$6\text{ kg}$}
{\True $4\text{ kg}$}
{$8\text{ kg}$}
\loigiai{
$2 + 2 = 4\text{ kg}$.\\
\textbf{Đáp án của câu hỏi là C}
}
\end{ex}

\begin{ex} %[2H3V17]
Rót hết nước từ can $10\text{ l}$ sang hai bình: bình một được $6\text{ l}$, bình hai được
\choice
{\True $4\text{ l}$}
{$3\text{ l}$}
{$5\text{ l}$}
{$6\text{ l}$}
\loigiai{
$10 - 6 = 4\text{ l}$.\\
\textbf{Đáp án của câu hỏi là A}
}
\end{ex}

\begin{ex} %[2H3V17]
Một bao gạo nặng $50\text{ kg}$, người ta chia đều vào 5 túi nhỏ. Mỗi túi nặng là
\choice
{$8\text{ kg}$}
{\True $10\text{ kg}$}
{$12\text{ kg}$}
{$15\text{ kg}$}
\loigiai{
$50 : 5 = 10\text{ kg}$.\\
\textbf{Đáp án của câu hỏi là B}
}
\end{ex}

\begin{ex} %[2H3C17]
Một con mèo và một con gà cùng đứng lên cân thì cân chỉ $5\text{ kg}$. Biết con mèo nặng hơn con gà $1\text{ kg}$. Con mèo cân nặng là
\choice
{\True $3\text{ kg}$}
{$2\text{ kg}$}
{$4\text{ kg}$}
{$1\text{ kg}$}
\loigiai{
Con mèo: $(5 + 1) : 2 = 3\text{ kg}$.\\
\textbf{Đáp án của câu hỏi là A}
}
\end{ex}

\textbf{Phần 2. Trắc nghiệm đúng sai}
\setcounterpageref{ex}{0}

\begin{ex} %[2H3H17]
Xét tính hợp lý trong thực tế:
\choiceTF
{\True Một quả trứng gà nặng khoảng vài chục gam (bé hơn $1\text{ kg}$).}
{\True Một con bò cân nặng vài trăm ki-lô-gam.}
{Một người trưởng thành uống khoảng $20\text{ l}$ nước mỗi ngày.}
{\True Mỗi ngày người trưởng thành uống khoảng $1\text{ l}$ đến $2\text{ l}$ nước.}
\loigiai{
\begin{itemize}
    \item Trứng gà bé hơn $1\text{ kg}$: Đúng.
    \item Con bò nặng vài trăm $\text{kg}$: Đúng.
    \item Uống $20\text{ l}$ là phi thực tế. Ý 3 Sai, Ý 4 Đúng.
\end{itemize}
\begin{itemize}
    \item a) \textbf{Đúng}.
    \item b) \textbf{Đúng}.
    \item c) \textbf{Sai}.
    \item d) \textbf{Đúng}.
\end{itemize}
}
\end{ex}

\begin{ex} %[2H3V17]
Có 2 can dầu: Can A chứa $15\text{ l}$, can B chứa $20\text{ l}$. Xét tính đúng sai:
\choiceTF
{\True Cả hai can chứa $35\text{ l}$ dầu.}
{\True Can B chứa nhiều hơn can A là $5\text{ l}$ dầu.}
{Nếu rót $5\text{ l}$ từ can B sang can A thì can B còn $10\text{ l}$.}
{\True Khi đó can B còn $20 - 5 = 15\text{ l}$ và hai can bằng nhau.}
\loigiai{
\begin{itemize}
    \item Tổng: $15 + 20 = 35\text{ l}$: Đúng.
    \item Hiệu: $20 - 15 = 5\text{ l}$: Đúng.
    \item $20 - 5 = 15\text{ l} \neq 10\text{ l}$. Ý 3 Sai, Ý 4 Đúng.
\end{itemize}
\begin{itemize}
    \item a) \textbf{Đúng}.
    \item b) \textbf{Đúng}.
    \item c) \textbf{Sai}.
    \item d) \textbf{Đúng}.
\end{itemize}
}
\end{ex}

\textbf{Phần 3. Trả lời ngắn}
\setcounterpageref{ex}{0}

\begin{ex} %[2H3H17]
Một bao xi măng nặng $50\text{ kg}$, một bao cát nặng $40\text{ kg}$. Tính tổng khối lượng cả hai bao theo ki-lô-gam.
\shortans{90}
\loigiai{
\textbf{Đáp số:} $90$\\
$50 + 40 = 90\text{ kg}$.
}
\end{ex}

\begin{ex} %[2H3H17]
Một thùng nước chứa $60\text{ l}$, đã dùng hết $25\text{ l}$. Còn lại bao nhiêu lít nước?
\shortans{35}
\loigiai{
\textbf{Đáp số:} $35$\\
$60 - 25 = 35\text{ l}$.
}
\end{ex}

\begin{ex} %[2H3V17]
Mẹ mua 3 túi táo, mỗi túi nặng $2\text{ kg}$. Tính tổng cân nặng của 3 túi táo theo ki-lô-gam.
\shortans{6}
\loigiai{
\textbf{Đáp số:} $6$\\
$2 \times 3 = 6\text{ kg}$.
}
\end{ex}

\begin{ex} %[2H3V17]
Có 4 can nước mắm, mỗi can chứa $5\text{ l}$. Tính tổng số lít nước mắm trong 4 can.
\shortans{20}
\loigiai{
\textbf{Đáp số:} $20$\\
$5 \times 4 = 20\text{ l}$.
}
\end{ex}

\begin{ex} %[2H3V17]
Bé An nặng $24\text{ kg}$, bé Bình nhẹ hơn bé An $3\text{ kg}$. Hỏi bé Bình nặng bao nhiêu ki-lô-gam?
\shortans{21}
\loigiai{
\textbf{Đáp số:} $21$\\
$24 - 3 = 21\text{ kg}$.
}
\end{ex}

\begin{ex} %[2H3C17]
Một thùng to đựng $30\text{ l}$ nước, rót đầy sang một can $12\text{ l}$ và một can $8\text{ l}$. Hỏi trong thùng to còn lại bao nhiêu lít nước?
\shortans{10}
\loigiai{
\textbf{Đáp số:} $10$\\
$30 - (12 + 8) = 10\text{ l}$.
}
\end{ex}
